\section{Research Questions and Evidence Map}\label{sec:evidence}
Each claimed contribution below is tied to a research question (RQ), a controlled experiment, a quantitative result with its uncertainty, a comparison, a limitation and a reproducible artifact (paths relative to the repository root; all values are generated from \texttt{results/final/final\_results.json}). Contributions that could not be supported are listed as negative or null results.

\begin{description}
\item[C1 -- Observable locality (RQ1: does a local-sum observable improve a matched QNN?).] \emph{Experiment:} \LsSeeds{} paired seeds with identical initialization per seed. \emph{Result:} mean PR-AUC difference \LsDelta{}, local-sum better in \LsWins{} seeds. \emph{Uncertainty:} bootstrap 95\% CI [\LsCIlo{}, \LsCIhi{}]; one-sided Wilcoxon $p=\LsWilcoxon{}$; permutation $p=\LsPerm{}$. \emph{Comparison:} global-observable QNN. \emph{Limitation:} controlled subset, one dataset and circuit family; architecture-level effect only. \emph{Artifact:} \texttt{results/layer1/local\_sum\_decisive/statistical\_tests.json}.
\item[C2 -- Boundary of the quantum benefit (RQ2: does C1 yield system-level value?).] \emph{Experiment:} locked confirmation split with a matched classical nonlinear control. \emph{Result:} XGBoost \CapXgb{}; $+$local-sum QNN \CapXgbLocal{}; $+$classical control \CapXgbControl{}. \emph{Uncertainty:} CI of the quantum increment [\CapCIlo{}, \CapCIhi{}] includes zero. \emph{Comparison:} XGBoost alone and the classical control. \emph{Limitation:} one confirmation split. \emph{Outcome:} no verified quantum predictive advantage (negative result). \emph{Artifact:} \texttt{results/layer1/local\_sum\_decisive/incremental\_value\_result.json}.
\item[C3 -- Implementation versus primitive cost in PQC (RQ3: how strongly does the backend change measured PQC latency?).] \emph{Experiment:} identical messages and parameter sets on the pure-Python reference and compiled \texttt{pqcrypto}; warm-ups excluded. \emph{Result:} speedups \SpeedupMin{}--\SpeedupMax{}$\times$ (ML-KEM-768 encapsulation \SpeedupMlkemEncaps{}$\times$; ML-DSA-65 signing \SpeedupMldsaSign{}$\times$; SLH-DSA-128s signing \SpeedupSlhSSign{}$\times$, compiled median \OptSlhSSignMs{}~ms). \emph{Uncertainty:} P5--P95 and 95\% $t$-intervals per operation; $n$ reported per backend (reference SLH-DSA signing $n=\SlhSRefN{}$). \emph{Comparison:} reference versus compiled; classical ECDH/ECDSA/RSA by role. \emph{Limitation:} not native \texttt{liboqs}; unequal repetition counts; machine-specific. \emph{Artifact:} \texttt{results/layer2/pqc\_backend\_comparison.csv}, \texttt{pqcrypto\_raw.csv}.
\item[C4 -- End-to-end cost of protection (RQ4: what does quantum-resistant protection cost per transaction?).] \emph{Experiment:} serialize--sign--key-establish--AES-256-GCM--decrypt--verify for six configurations. \emph{Result:} median \EEAMedian{}~ms (classical), \EECMedian{}~ms (compiled ML-KEM-768 + ML-DSA-65), \EEBMedian{}~ms (reference), \EEDMedian{}~ms (hybrid); payload expansion \EEAExpansion{}$\times$ versus \EECExpansion{}$\times$. \emph{Uncertainty:} P95 and $n$ per configuration. \emph{Comparison:} classical baseline. \emph{Limitation:} in-memory transmission, pre-distributed keys. \emph{Artifact:} \texttt{results/layer2/end\_to\_end\_banking\_summary.csv}.
\item[C5 -- Independent BB84 cross-validation (RQ5: do a protocol simulator and explicit circuits agree with theory and each other?).] \emph{Experiment:} \ProtoTrials{} protocol trials and \QiskTrials{} Qiskit trials per condition across six conditions. \emph{Result:} intercept--resend mean QBER \BbqProtoKEve{} (protocol) and \BbqQiskEve{} (Qiskit) against theory 0.25; MAE \BbMAE{}, RMSE \BbRMSE{}, maximum gap \BbMaxDiff{}. \emph{Uncertainty:} 95\% $t$-intervals (e.g.\ \BbciProtoKEve{} and \BbciQiskEve{}). \emph{Comparison:} theory and the other implementation. \emph{Limitation:} different noise models; 500 bits for Qiskit; software simulation only. \emph{Artifact:} \texttt{results/layer2/bb84\_final\_summary.csv}, \texttt{bb84\_final\_agreement.json}.
\item[C6 -- Eve-detection operating point (RQ6: is 0.11 a defensible threshold?).] \emph{Experiment:} threshold sweep over Eve-absent and Eve-present trials. \emph{Result:} at 0.11, TPR \ThrTPRProtoK{} and FPR \ThrFPRProtoK{} (protocol, 2000 bits) when 10\%-noise trials count as benign. \emph{Uncertainty:} three data sets and two noise ranges. \emph{Comparison:} balanced-accuracy optimum. \emph{Limitation:} optimum depends on the noise range; QBER does not prove eavesdropping. \emph{Artifact:} \texttt{results/layer2/bb84\_threshold\_sweep.csv}.
\item[C7 -- Simulator-to-hardware validation (RQ7: how much error do real devices add, and does the attack stay distinguishable?).] \emph{Experiment:} \HwJobs{} independent \HwBackend{} jobs, pre-selected best/typical/worse qubits, ideal and backend-derived noisy Aer controls on identical mappings. \emph{Result:} no-Eve QBER \HwNoEveBest{}, \HwNoEveTypical{}, \HwNoEveWorse{}; intercept--resend \HwEveBest{}, \HwEveTypical{}, \HwEveWorse{}; Eve--no-Eve separation in all 15 job-by-qubit cells (\HwMaxNoEve{} $<$ \HwMinEve{}). \emph{Uncertainty:} 95\% intervals over jobs (best qubit no-Eve \HwNoEveCIBest{}); bootstrap CI of the attack increment [\HwDeltaBootLoBest{}, \HwDeltaBootHiBest{}]; with $n=5$ the smallest attainable permutation $p$ is \HwPermBest{}. \emph{Comparison:} theory, ideal Aer, noisy Aer (MAE \HwNoisyMAE{} versus ideal-Aer MAE \HwIdealMAE{}). \emph{Limitation:} one backend, one calibration window, three qubits, circuit-level only (not physical QKD). \emph{Artifact:} \texttt{results/ibm\_quantum/hardware\_statistics.json}, raw counts under \texttt{backend\_runs/A}.
\item[C8 -- Executability of the QML circuits (RQ8: do the Layer~1 circuits run on hardware, and how different are the outputs?).] \emph{Experiment:} \texttt{ibm\_fez}, frozen simulator-trained weights, six held-out transactions. \emph{Result:} mean absolute deviation from simulation \FezCltMAD{} (Chow--Liu) and \FezLinearMAD{} (linear). \emph{Uncertainty:} none estimable (single job, six samples). \emph{Comparison:} simulator expectation values. \emph{Limitation:} executability only; no hardware training. \emph{Artifact:} \texttt{reports/ibm\_hardware/hardware\_run.json}.
\end{description}
